\subsection{Design Principles}
\label{sec:principle}
\begin{AIbox}
The design of GenericAgent (GA) is guided by two principles: \textbf{high contextual information density} and \textbf{continual self-evolution}. The first principle targets a key limitation in LLM-based agents, where under a fixed context budget, performance depends less on context length and more on how much decision-relevant information each token carries. The second principle enables the system to improve over time through accumulated experience, gradually extending its capabilities and shifting the agent from a static prompt-driven setup to a continuously adapting system.
\end{AIbox}


\subsubsection{Maximizing Contextual Information Density: From Long Context to High-Value Context}

The central objective of this principle is to maximize decision-relevant information under a fixed context budget. In practice, performance degradation in LLM-based agents rarely comes from insufficient context length. It is more often caused by inefficient context organization. Redundant information dilutes attention, irrelevant historical memory interferes with current reasoning, and overly compressed representations obscure semantic structure and reduce reliability. These failure modes share a common cause: low information density. To address this, GA requires that all contextual content satisfy three constraints:
\begin{itemize}
    \item \textbf{Completeness.} {\small
    All information required for the current decision must be explicitly present in the context, preventing the model from relying on implicit assumptions or hallucinated completions.
    }

    \item \textbf{Conciseness.} {\small
    Irrelevant and redundant information must be eliminated to reduce noise and ensure that attention is focused on decision-critical signals.
    }

    \item \textbf{Naturalness.} {\small
    The context must be expressed in well-formed natural language that preserves semantic clarity, avoiding overly compressed or artificial representations that hinder comprehension.
    }
\end{itemize}

These three constraints form a minimal and complementary set. Completeness ensures that no required information is missing. Conciseness removes all non-essential content. Naturalness preserves the interpretability of the remaining information. Together, they define a high-density contextual representation that is both information-complete and cognitively tractable.

Based on these constraints, GA introduces four system-level mechanisms to enforce high-density context construction. (1) First, it adopts a minimal and orthogonal tool interface composed of a fixed set of atomic tools. This prevents tool-space expansion and removes redundant interface definitions, improving conciseness at the action level. (2) Second, GA employs an index-driven hierarchical memory mechanism.  Historical information is first filtered through a lightweight index and only selectively expanded when needed. This ensures completeness without requiring full-history injection. (3) Third, repeated workflows are distilled into structured SOPs and compiled into executable skill units. This replaces multi-step reasoning traces with compact reusable abstractions, improving conciseness while preserving structural clarity. (4) Finally, web content is processed through DOM simplification and incremental updates. Only semantically meaningful content and changes are introduced into the context, which maintains both conciseness and naturalness. 

Through these mechanisms, GA shifts context construction from passive accumulation to active optimization. The context is transformed from a long sequence of tokens into a compact, high-value information substrate for decision making.

\subsubsection{Continual Self-Evolution: From One-Shot Prompt Design to Runtime Skill Growth}

This principle describes how GA evolves over time. GA is not a fixed prompt wrapped around a model. It is a agent system whose internal behavior is continuously shaped by execution, while its external interface remains stable and minimal.

The core mechanism is execution-driven self-evolution. During task execution, GA transforms raw execution traces into structured representations of its own decisions, rather than leaving them as unorganized dialogue history. As similar tasks recur, these representations are progressively distilled and compressed into reusable strategies and executable skills, enabling future problem solving to rely less on repeated ad-hoc reasoning and more on internal skill composition.

\subsubsection{A Self-Reinforcing Loop Between Two Principles}

GA operates as a unified agent system where contextual information density and continual self-evolution form a mutually reinforcing loop. High-density contexts ground execution in precise, decision-relevant signals. This makes experience from runtime behavior more structured and reliable. These experiences are then compressed into SOPs and reusable skills, reducing reliance on repeated token-level reasoning and enabling more direct skill-based execution in future tasks.

Over time, the two processes reinforce each other. Better structured contexts improve experience abstraction, while accumulated skills reduce the need for extensive reasoning. This allows future executions to operate with more compact inputs. GA thus evolves by improving both how information is represented and how it is turned into behavior under repeated problem settings.



\subsection{Overview of GenericAgent}
\label{sec:overview}

GA is a general-purpose LLM agent system that provides automation over a local computing environment. It enables a compatible LLM backend to operate across browsers, terminals, file systems, keyboard and mouse inputs, screen-based perception, and mobile devices through configurable tool adapters. The system is model-agnostic, allowing the underlying reasoning engine (e.g., Claude, GPT, Gemini) to be replaced without affecting execution logic, tool interfaces, or memory architecture. The overall architecture is illustrated in Figure \ref{fig:agent_loop}.

GA is organized around a unified agent loop that connects the model and the environment. At each step, the system combines global memory with the current task specification to form an execution context. The LLM processes this context and produces either direct outputs or tool calls. Tool execution is delegated to external modules, and the results are returned as structured signals that update the system state. This loop defines a continuous interaction between reasoning and environment feedback, enabling tool-grounded iterative execution.

On top of this loop, GA supports two execution modes. \textbf{Interactive execution} handles user-initiated tasks such as code execution, information retrieval, and file operations. \textbf{Reflective execution} runs in the background and is triggered by schedules or system events. It is responsible for monitoring, memory consolidation, and experience distillation. Both modes share the same agent loop and memory system, differing only in initiation and output handling. To support long-horizon operation, GA defines explicit execution bounds that ensure controllability while preserving interruptibility and scope constraints.

After task completion, GA performs memory distillation. Execution traces are compressed into structured long-term representations and stored in shared memory. These representations are reused across both execution modes, allowing accumulated experience to improve future execution through compression and reuse rather than unbounded computation.

\subsection{Core Components of GA}
\label{sec:core_design}
Building on the principles in Section \ref{sec:principle}, we describe the core components of GA, covering the atomic tool set, hierarchical memory management, reflection-driven self-evolution, and structured-extraction-based browser interaction.

\subsubsection{Minimal Atomic Tool Set}
\label{sec:tools}
In existing agent loop architecture, tool definitions are explicitly injected into the context at every LLM call, making tool design a key determinant of context information density. To reduce the overhead introduced by tool schemas, GA adopts a minimal tool set consisting of nine atomic tools and follows two core principles: \textbf{atomicity} and \textbf{compositional generalization}. Atomicity constrains tool granularity by enforcing irreducible primitive capabilities, while compositional generalization enables the expression of complex tasks through sequential composition of these atomic tools. Table~\ref{tab:tools} summarizes the complete GA tool set, which covers core capabilities including code execution, file manipulation, web content retrieval and browser-side execution, explicit user interaction, and triggers for working-memory updates and long-term memory refinement.

This minimal tool set yields two direct benefits. First, it significantly reduces tool-description overhead in the context. In multi-turn interactions, a smaller and more interface-efficient tool set directly reduces cumulative token consumption across calls. Second, it reduces the model’s decision complexity. A compact and well-bounded action space decreases uncertainty in tool selection, enabling more stable tool-use patterns and thereby reducing execution errors and retries.

Besides, GA does not rely on specialized tools to expand functionality. Instead, complex behaviors are composed from atomic primitives. For example, file-level operations follow a standard read–modify–write paradigm, while web interactions are decomposed into content retrieval and browser-side execution. External capabilities are encapsulated as Skills and invoked indirectly through general-purpose code-execution and web-related primitives, and are loaded into the context only on demand. 


\begin{table}[t]
    \centering
    \caption{The nine atomic tools in GenericAgent.}
    \label{tab:tools}
    {\footnotesize
    \setlength{\tabcolsep}{3pt}
    \renewcommand{\arraystretch}{0.95}
    \begin{tabularx}{\linewidth}{@{}lX@{}}
    \toprule
    \textbf{Tool} & \textbf{Brief Description} \\
    \midrule
    \texttt{code\_run} &
    Execute Python or shell code with timeout and abort control. \\
    \texttt{file\_read} &
    Read local files by line range, with optional keyword-based location. \\
    \texttt{file\_write} &
    Create new files or write large content blocks (overwrite/append/prepend). \\
    \texttt{file\_patch} &
    Apply precise, context-checked patch updates to existing file regions. \\
    \texttt{web\_scan} &
    Extract the current browser page as simplified HTML or text, with tab metadata. \\
    \texttt{web\_execute\_js} &
    Inject and execute JavaScript in the user’s browser, capturing return values. \\
    \texttt{ask\_user} &
    Issue a synchronous query to the user and wait for an explicit response. \\
    \shortstack[l]{\texttt{update\_working\_}\\\texttt{checkpoint}} &
    Update a short-term working notepad injected in subsequent reasoning turns. \\
    \shortstack[l]{\texttt{start\_long\_term\_}\\\texttt{update}} &
    Trigger long-term memory distillation over recent execution traces. \\
    \bottomrule
    \end{tabularx}
    }
    \end{table}

\subsubsection{Hierarchical Memory Mechanism}
\label{sec:momory}

GA retains only information that cannot be easily derived from the current context or the model’s own reasoning. The goal is to maximize behavioral impact per token under a limited context budget. The memory system follows a four-layer architecture with hierarchical distillation and taxonomy-based organization.

\textbf{L0: Meta-Rule Layer.}  L0 defines the core behavioral constraints of GA. It acts as a lightweight control layer for system operation and self-regulation. Each rule is written in a single sentence. For example: ``Only information verified through execution can be written to long-term memory.'' This rule prevent unverified or hallucinated information from entering long-term memory. L0 is injected into the system prompt at startup with minimal overhead and remains fixed during execution.

\textbf{L1: Index Layer.} L1 provides a lightweight index for memory routing. It stores only metadata, file paths, short descriptions, and applicability conditions. It does not store full content.

Before each task, GA queries L1 and retrieves only relevant memory entries. This keeps the active context sparse even as the memory base grows. The L1 index is maintained at around 30 lines, keeping retrieval overhead low.

\textbf{L2: Fact Store.}  L2 stores verified environmental facts. It includes information that cannot be easily inferred, such as directory structures, configurations, API endpoints, and system constraints.

All entries are validated through tool execution or external signals. Information is organized by topic and incrementally summarized. This reduces fragmentation and improves retrieval clarity.

\textbf{L3: Task-Level SOP Layer.} L3 stores reusable task-level procedures. It captures patterns such as tool usage order, dependencies, failure cases, and correction strategies.

These SOPs are extracted from execution traces after task completion. They are not manually written. In later tasks, a small amount of L3 content can guide behavior and improve success rates.

In summary, the memory system does not aim to store as much information as possible. It focuses on keeping only behaviorally useful knowledge. L0 defines constraints. L1 handles routing. L2 stores verified facts. L3 captures reusable experience. Together, they allow GA to operate within a fixed context budget while maximizing the utility of each token.

\subsubsection{Reflection-Driven Self-Evolution Mechanism}
\label{sec:self-evolution}

GA’s self-evolution mechanism is an abstraction-driven process that compiles execution trajectories into high-density reusable representations. It preserves semantic integrity while reducing reasoning cost. Through periodic reflection, the system analyzes historical execution traces and converts them into structured knowledge. These structures are then organized into executable forms. The evolution is guided by user task distributions and behavioral patterns, ensuring alignment with real-world usage and observed failures. The mechanism consists of three stages.

\textbf{Stage 1: Experience Distillation into SOPs.} GA analyzes historical execution traces to extract useful information, including successful decision paths and failure patterns. These are abstracted into SOPs that describe stable execution logic.

This stage serves two purposes. First, it generalizes instance-level behavior into reusable process specifications. Second, it prioritizes high-frequency and high-error tasks. Routine or low-impact steps are discarded to improve compression and generalization.

\textbf{Stage 2: SOP Compilation into Code.} SOPs with stable structure are evaluated for compilation. If a task shows clear parameterization, deterministic logic, and explicit dependencies, it is converted into code.

This replaces procedural reasoning with program execution. Execution becomes structured and deterministic. Compared to SOPs, code provides clearer semantics and more stable behavior under repeated execution.

\textbf{Stage 3: Registration as Skills.} Compiled code modules are registered as Skills in GA’s indexing system. Skills are reusable execution units. Similar tasks can directly invoke them, without intermediate reasoning.

GA monitors Skill usage and updates implementations over time. Updates include interface simplification, parameter refinement, and execution optimization. This improves efficiency per token.

These stages form a closed loop. Reflection extracts experience, SOPs generalize behavior, code replaces reasoning, and Skills enable reuse. This process improves efficiency over time and reduces redundant computation. 

\subsubsection{Efficient Browser Interaction via Structured Extraction}
GA treats browser interaction as structured extraction of GA treats browser interaction as structured extraction of task-relevant information, instead of full-page retrieval. This avoids exposing verbose HTML, which contains large amounts of irrelevant markup and leads to low information density.

At the content layer, GA introduces simpHTML, which converts web pages into compact structured representations. It parses the DOM, removes non-informative regions, and identifies repeated content patterns. Candidate regions are ranked using geometric and structural signals, and only the most informative content is retained. The result is a compressed Markdown representation that reduces token usage while preserving task-relevant information.

At the execution layer, GA uses a Chrome Extension with direct access to the Chrome DevTools Protocol (CDP), instead of higher-level automation frameworks. This enables direct interaction with complex page structures, including cross-origin iframes, Shadow DOM, and multi-tab environments. Combined with DOM traversal and coordinate transformation, GA can handle real-world pages more reliably than scripted automation.

Overall, GA reduces browser interaction from full-page retrieval to structured extraction, lowering token cost and improving interaction robustness.







